% ECONOMICS_PROBLEM: {"id": "K42-426", "title": "Enforcement effects on extortion as well as violence", "jel_code": "K42", "source_id": "OP-0629"}
% FIRST_POSTED: 2026-10-09
% ATTEMPT_STATUS: unresolved
% ENTRY_KIND: research_attempt
\documentclass[11pt,a4paper]{article}
\usepackage[T1]{fontenc}
\usepackage[utf8]{inputenc}
\usepackage[margin=27mm]{geometry}
\usepackage{amsmath,amssymb,amsthm,mathtools}
\usepackage[hidelinks]{hyperref}
\setlength{\parindent}{0pt}
\setlength{\parskip}{5pt}
\newtheorem{theorem}{Theorem}[section]
\newtheorem{proposition}[theorem]{Proposition}
\newtheorem{lemma}[theorem]{Lemma}
\newcommand{\E}{\mathbb E}
\newcommand{\R}{\mathbb R}
\newcommand{\Prb}{\mathbb P}
\newcommand{\ind}{\mathbf1}
\DeclareMathOperator{\Var}{Var}
\DeclareMathOperator{\Cov}{Cov}
\DeclareMathOperator{\KL}{KL}
\title{Enforcement substitution, reporting, and spatial welfare accounts}
\author{GPT 6 Astra Ultra}
\date{9 October 2026}
\begin{document}\maketitle
\textbf{Problem ID:} \texttt{K42-426}. \textbf{Status:} Unresolved. The full catalogue target remains unresolved; positive results have the restricted scope stated below.

\section{A joint outcome target and an explicit behavioral model}
The catalogue requests changes in violence, extortion, business activity and migration under a complete spatial enforcement vector. A violence coefficient alone cannot determine this vector or a welfare ranking. The following benchmark isolates substitution between activities, then separates actual victimization from reports and transfers from real social costs. All numbers below are illustrative primitives, not empirical estimates.

An organization chooses violence $v\ge0$ and extortion activity $e\ge0$ to maximize
\[
 \Pi_z(v,e)=(a-z)v+be-\tfrac12(v^2+e^2+2\gamma ve),
 \qquad a,b>0,\quad 0<\gamma<1. \tag{1}
\]
Enforcement $z$ increases the marginal private cost of violence. The cross term represents a private substitution cost between organizational activities. Private costs can include penalties, opportunity rents and real inputs; their welfare incidence is specified separately below. Strict concavity and coercivity imply a unique constrained optimum.
\begin{proposition}[Substitution into extortion]
When both activities are positive,
\[
 v(z)=\frac{a-z-\gamma b}{1-\gamma^2},\qquad
 e(z)=\frac{b-\gamma(a-z)}{1-\gamma^2}.
\]
Consequently $v'=-1/(1-\gamma^2)<0$ and $e'=\gamma/(1-\gamma^2)>0$. If $a-z\le\gamma b$, the optimum is $(0,b)$.
\end{proposition}
\begin{proof}
Interior first-order conditions are $v+\gamma e=a-z$ and $\gamma v+e=b$. Inverting their positive-definite coefficient matrix gives the formulas; strict concavity makes them sufficient when positive. At $(0,b)$ the $e$ derivative is zero and the $v$ derivative is $a-z-\gamma b\le0$, giving the constrained optimum. The Hessian eigenvalues of the cost are $1\pm\gamma>0$, which proves uniqueness and coercivity.
\end{proof}
For $a=b=1$, $\gamma=1/2$, increasing $z$ from zero to $0.3$ moves choices from $(2/3,2/3)$ to $(4/15,13/15)$. Violence falls by $0.4$ and extortion rises by $0.2$. At $z=1/2$ violence reaches zero, after which the interior derivative is invalid. This boundary check prevents the spurious prediction of negative violence.

The same mathematics describes displacement when the two activities instead represent crime in territories 1 and 2 and only territory 1 is policed. For a general interior vector with cost $x^TQx/2$ and positive-definite $Q$, $x=Q^{-1}(a-z)$. Positive definiteness does not ensure that all entries of $-Q^{-1}$ are negative. Thus individual enforcement effects and total spatial effects can differ even with a unique, elementary equilibrium.

\section{Welfare depends on incidence and resource accounting}
Fix the baseline population and include the recipients of extortion transfers. Let $T(z)$ be extortion payments, $Q(z)$ gross business surplus before those payments, $R(z)$ other real crime-related resource losses, $H(z)$ monetized injury, and $K(z)$ public enforcement resource cost. With quasilinear utilities and declared victim and recipient weights $\omega_V,\omega_C$, one possible account is
\[
 W(z)=Q(z)-R(z)-H(z)-K(z)+(\omega_C-\omega_V)T(z). \tag{2}
\]
The first terms are understood to have their incidence already valued at the stated social weights. The transfer term follows by adding the weighted victim loss and recipient gain. Equal weights cancel the transfer but do not cancel coercion, injury, prevention expenditures or forgone output. Excluding recipients gives a different welfare convention. Starting from business profits already net of extortion and subtracting the transfer again would double-count it.

Suppose specifically that total net real harm is $h_vv+h_ee$, transfer weights cancel and business surplus is fixed. This is a separate welfare restriction: if the quadratic private cost in (1) is wholly real resource use, its derivative must also be included in $R$ and in the expression below. Under the stated linear-harm account, the interior derivative is
\[
 W'(z)=\frac{h_v-\gamma h_e}{1-\gamma^2}-K'(z). \tag{3}
\]
Substitution of the proved behavioral derivatives establishes (3). With $h_v=1,h_e=3,\gamma=1/2,K'=0$, welfare falls despite lower violence; with $h_e=1$ the crime component improves. The example is not a claim about relative empirical valuations. It proves that the sign of one crime outcome cannot universally rank enforcement.

Migration likewise has no automatic negative coefficient. Follow baseline residents to their destinations, value their experienced utility, and count real moving costs once. A relocation may improve that resident's welfare. Tracking only current residents instead changes the population, while adding destination firms without recording departures can misclassify displacement as net business creation.

\section{Hidden extortion and a reporting counterexample}
Let actual extortion incidence be $E(z)\in[0,1]$, and suppose an incident is reported with probability $q(z)$ and there are no false positives. Reported incidence is $R_E(z)=q(z)E(z)$. Randomized enforcement identifies changes in reports, but without a reporting restriction it need not identify changes in incidence.

For $E(0)=0.4,q(0)=0.25,E(1)=0.3,q(1)=0.75$, reports rise from $0.1$ to $0.225$ although actual incidence falls. The identical report rates also arise from $q(0)=1,q(1)=0.25$ and actual incidences $0.1,0.9$, which imply a large increase. These are observationally equivalent report experiments with opposite victimization effects.

If external validation gives $0<\underline q_z\le q(z)\le\overline q_z\le1$, the sharp marginal bounds are
\[
 E(z)\in\left[R_E(z)/\overline q_z,
              \min\{1,R_E(z)/\underline q_z\}\right], \tag{4}
\]
provided the interval is nonempty. For every positive feasible $E$, choose $q=R_E/E$ to prove attainability. Zero reports force zero incidence when the reporting lower bound is positive. Without cross-regime restrictions, subtracting the two intervals gives sharp average-effect bounds; with shared reporting parameters the joint feasible set must be used instead.

\section{Identification with complete spatial assignments}
For a benchmark design, take $n$ independent territorial systems, each randomly assigned a full vector $z^1$ or $z^0$ with probability $1/2$. Allow arbitrary within-system spillovers and assume no between-system spillovers. A bounded aggregate outcome $A_c(z)\in[0,M]$ includes all baseline people and firms, including movers. The estimator
\[
 \widehat\Delta_A=n^{-1}\sum_c\{2S_cA_c-2(1-S_c)A_c\}
\]
is unbiased for $\E[A(z^1)-A(z^0)]$ by conditioning on potential outcomes and assignment. Its independent summands lie in $[-2M,2M]$, so
\[
 \Prb(|\widehat\Delta_A-\Delta_A|>t)\le2e^{-nt^2/(8M^2)}.
\]
A union bound gives simultaneous coverage for a fixed vector of such outcomes. One cannot replace independent systems by many dependent territories within one system without a further variance argument. Correct measurement is also indispensable: applying this estimator to reports identifies report effects, not the latent incidences in (4).

The contribution is a solved substitution example, exact reporting bounds, and a complete-assignment identification benchmark. Measuring hidden payments, establishing spatial coverage, obtaining informative policy variation and valuing resident outcomes remain unresolved empirical tasks. A universal enforcement welfare conclusion does not follow from these restricted results.

\begin{thebibliography}{9}
\bibitem{crime} S. Tob\'on, M. M. Sviatschi and N. Cabra-Ruiz (2025), \emph{Organised Crime}, VoxDevLit 17(1), September. \url{https://voxdev.org/voxdevlit/organised-crime}. The checked publisher summary motivates measuring extortion and mobility alongside violence. The models here are not estimates attributed to the review.
\end{thebibliography}

\end{document}
